% Inside the appendices, \appendix takes the title where \chapter would.
\appendix{Viết tắt và thuật ngữ chính\label{appendix:terms}}
Các viết tắt và thuật ngữ dùng xuyên suốt luận văn, cho độc giả ngoài lĩnh vực viễn thông.

\begin{table}[H]
\centering
\caption[Viết tắt và thuật ngữ chính]{Viết tắt và thuật ngữ chính dùng xuyên suốt luận văn.\label{tab:abbrev}}
% One size down so that a glossary of this length shares a page with the appendix
% heading inside the class's 23cm text block. The \singlespacing is redundant while
% the body is single-spaced, and is kept so the table survives a change back.
\singlespacing
\footnotesize
\renewcommand{\arraystretch}{1.2}
\begin{tabular}{@{}p{0.21\linewidth} p{0.71\linewidth}@{}}
\hline\hline
Thuật ngữ & Ý nghĩa \\
\hline
RAN & Mạng truy nhập vô tuyến: phần mạng di động nối thiết bị với lõi \\
NG-RAN & RAN thế hệ tiếp theo, kiến trúc mạng truy nhập vô tuyến 5G \\
IMT-2030 & Khung ITU cho thế hệ sau 5G, thường gọi là 6G \\
AI-RAN & Kiến trúc trong đó AI và khối lượng RAN chia sẻ một hạ tầng tăng tốc \\
gNB & Trạm gốc 5G (``next-generation NodeB'') \\
gNB-CU / DU & Đơn vị trung tâm / Đơn vị phân tán: hai nửa tách rời của gNB \\
gNB-CU-CP & Phần mặt điều khiển của đơn vị trung tâm: lớp hệ thống được nghiên cứu ở đây \\
UE & Thiết bị người dùng: thiết bị đầu cuối (điện thoại, modem) gắn vào mạng \\
UECM & Quản lý ngữ cảnh UE: phân hệ quản lý trạng thái theo UE được hiện thực lại ở đây \\
RRC & Điều khiển tài nguyên vô tuyến: giao thức báo hiệu mặt điều khiển giữa UE và gNB \\
DRB & Kênh mang dữ liệu vô tuyến theo UE: kênh dữ liệu mặt người dùng thiết lập trong phiên \\
Mặt điều khiển & Đường báo hiệu và quản lý, trái với đường dữ liệu mặt người dùng \\
Độ trễ đuôi & Độ trễ phần trăm cao (p99, p99.9), chỉ số thời gian hướng trường hợp xấu nhất \\
Jitter & Mức độ độ trễ dao động giữa các thao tác, trái với giá trị điển hình \\
Bão hoà & Tải đề nghị vượt quá khả năng dịch vụ theo kịp và yêu cầu bắt đầu xếp hàng \\
Shard & Một trong các lát trạng thái UE độc lập, mỗi lát thuộc đúng một luồng worker \\
FFI & Giao diện hàm nước ngoài: gọi Rust từ C++ qua ABI C \\
HDR histogram & Biểu đồ tần suất dải động cao, dùng để ghi độ trễ chính xác \\
\hline
\end{tabular}
\end{table}

\appendix{Giao thức tìm kiếm tài liệu\label{appendix:search}}
Mục~\ref{sec:approach} báo cáo hai phát hiện mà tìm kiếm này tạo ra. Giao thức được ghi ở đây để
khoảng trống mà nghiên cứu tuyên bố có thể được kiểm tra chứ không phải tin theo lời nói.

Tài liệu ứng viên được thu thập từ ACM Digital Library, IEEE Xplore, USENIX và Google
Scholar, dùng các tổ hợp thuật ngữ \emph{Rust}, \emph{C++}, \emph{memory safety}, \emph{tail
latency}, \emph{systems programming}, \emph{performance comparison}, và \emph{RAN} / \emph{5G
control plane}. Tìm kiếm giới hạn ở hội nghị có phản biện, trừ khi báo cáo ngành hoặc
chính phủ chính là đối tượng nghiên cứu, như các quy định an toàn bộ nhớ và thống kê lỗ hổng của
nhà cung cấp.
Tập kết quả sau đó được mở rộng bằng truy vết trích dẫn ngược và xuôi từ các bài an toàn Rust
cốt lõi, và lọc giữ công trình hoặc đo ngôn ngữ hệ thống bằng thực nghiệm, hoặc hình thức hoá bảo
đảm của Rust, hoặc báo cáo kinh nghiệm áp dụng. Tài liệu miền (đặc tả 3GPP và
khảo sát O-RAN) được thêm riêng để cố định bối cảnh kiến trúc.

Chỉ các nguồn trên sẽ bỏ sót công trình tiền nhiệm gần nhất, đó là lý do hai nguồn nữa được thêm.
Các luận văn do nhà cung cấp lưu trữ gần nghiên cứu này nhất (Mục~\ref{sec:relatedwork}) không
xuất hiện trong cơ sở dữ liệu trích dẫn nào: không có trong Crossref, OpenAlex, DBLP hay Semantic
Scholar. Truy vết xuôi hay ngược không chạm tới chúng, vì không có mục được lập chỉ mục trích dẫn
chúng và danh sách tài liệu tham khảo của chúng không được nộp. Chúng được tìm bằng cách tìm trực
tiếp trên kho lưu trữ thể chế, DiVA cho các trường Bắc Âu và Finna cho các trường Phần Lan. Cùng hai
tìm kiếm đó thiết lập sự vắng mặt ở phía kia. Finna, liên kết các kho luận văn đại học Phần Lan,
không trả luận văn nào về Rust trong bối cảnh truy nhập vô tuyến hay viễn thông. Tìm theo tác giả
thuộc Nokia trong các ấn phẩm được lập chỉ mục trên OpenAlex cũng không trả kết quả. Các thực thể của
hãng giữ khoảng ba vạn công trình được lập chỉ mục, vài trăm trong số đó về mạng truy nhập vô tuyến,
và một công trình nhắc Rust, trong một bản thảo trước in không liên quan về giám sát thực thi tin
cậy.

Một ghi chú thực tiễn thuộc cùng giao thức, vì nó làm méo tìm kiếm ngây thơ. Truy vấn văn bản tự do
cho \emph{Rust} trong các chỉ mục chung xếp bệnh học thực vật và hóa học ăn mòn trên ngôn ngữ lập
trình, nên cần truy vấn giới hạn trường hoặc lĩnh vực.

\appendix{Thiết kế thống kê: thành viên họ và khả năng phân giải\label{appendix:stats}}
Mục~\ref{sec:stats} cố định so sánh nào vào họ đã hiệu chỉnh và thiết kế nhìn sâu bao xa vào một
khác biệt. Ba điểm đằng sau các cam kết đó được ghi ở đây thay vì trong văn bản chính.

\paragraph{Vì sao một so sánh không phân giải được không miễn phí.}
Kiểm thử thông lượng đạt được dưới bão hoà sẽ đo lại điểm vận hành chứ không phải mã. Nó cũng tốn
các so sánh có thể phân giải điều gì, vì Holm chia ngân sách trên cả họ và một kiểm định không tách
gì vẫn chiếm một phần. Qua bão hoà lập luận tương tự chạy ngược. Điều mỗi biến thể vẫn duy trì ở đó
không bị ghim vào gì, nên loại trừ nó sẽ bỏ đi đại lượng quan sát duy nhất chế độ đó cung cấp.

\paragraph{Vì sao không có phần trăm nào được kiểm định khi quá tải.}
Dưới lịch gửi cố định dự định, một producer không theo kịp sẽ tụt lại tích luỹ, nên phần trăm đo ở
đó báo cáo tồn đọng tích luỹ trước cửa sổ cắt nó. Gắn giá trị $p$ vào cận dưới sẽ kiểm định bộ
điều khiển chứ không phải phần mềm.

\paragraph{Các sàn mà thiết kế có thể chạm.}
Một kiểm định chính xác trên $n$ cặp không thể trả $p$ hai phía dưới $2/2^{n}$, và sau hiệu chỉnh một
họ $m$ so sánh giá trị điều chỉnh nhỏ nhất có thể là $m$ lần số đó. Đó là nguồn gốc của
$\MinAttainableP$ và $\CorrectedFloor$ ở Mục~\ref{sec:stats}. Một chiến dịch có sàn đã hiệu chỉnh
nằm trên ngưỡng đang áp dụng không thể chứng minh ý nghĩa dù hiệu ứng lớn đến đâu, nên báo ``không
có ý nghĩa'' từ chiến dịch như vậy sẽ nói nhiều về số lần lặp hơn về ngôn ngữ.

\appendix{Cấu hình máy chủ và chiến dịch\label{appendix:config}}
Bảng~\ref{tab:testbed} liệt kê máy đo của Mục~\ref{sec:testbed} so với giao thức bare-metal mà số
đuôi cấp testbed yêu cầu. Bảng~\ref{tab:params} gom cấu hình dịch vụ, khối lượng công việc và
chiến dịch, các giá trị suy ra của chúng Mục~\ref{sec:workload} trích dẫn.

\begin{table}[H]
\centering
\small
\caption[Máy chủ đã dùng và giao thức testbed còn phải thực hiện]{Máy chủ mà kết quả được báo cáo đo trên, so với giao thức mà số đuôi cấp testbed yêu
cầu (\texttt{bench/TESTBED.md}). Các hàng hai cột trùng là thuộc tính của mã nguồn và
trình biên dịch. Các hàng khác nhau là điều một chiến dịch bare-metal vẫn mang lại. Chỉ khác biệt
cuối giờ còn tốn một kết quả, vì điểm gối bão hoà và so sánh tài nguyên đã được giải quyết trên máy
này: phần còn lại là đuôi xa dưới bão hoà.\label{tab:testbed}}
\renewcommand{\arraystretch}{1.2}
\begin{tabular}{@{}R{0.22\linewidth} R{0.35\linewidth} R{0.35\linewidth}@{}}
\hline\hline
Khía cạnh & Dùng cho kết quả báo cáo ở đây & Yêu cầu cho số đuôi cấp testbed \\
\hline
Máy chủ & Máy chủ hai socket chuyên dụng, $2\times$ Xeon E5-2680~v4, $16$ vCPU, một nút NUMA & Cùng lớp máy, không hypervisor \\
Ảo hoá & Khách KVM: tần số và turbo do host điều khiển và không quan sát được & Không. Governor và turbo đặt và xác minh \\
CPU steal & Khác không, cỡ $10$~ms thời gian vCPU mỗi giây & Không \\
Kernel & Rocky Linux~8.10, kernel bản phân phối & Giống, với \texttt{isolcpus} và \texttt{nohz\_full} trên lõi benchmark \\
Toolchain & \texttt{g++~15.1.1}, \texttt{rustc~1.97} & \texttt{g++~15.1.1}, \texttt{rustc~1.97} \\
Ghim CPU & Lõi $2$--$15$ qua \texttt{taskset}, $7$ shard, $7$ producer & Giống, trên lõi cũng bị gỡ khỏi bộ lập lịch \\
Đa luồng đồng thời & Vắng trong topology khách thấy. vCPU có chia anh em luồng hay không do host quyết & Vắng hoặc tắt, và xác minh trên máy \\
Khởi động nóng & Phần đầu mỗi lần chạy bỏ & Phần đầu mỗi lần chạy bỏ \\
Lần lặp và xen kẽ & \NRepeats{} lần lặp, bốn biến thể xen kẽ mỗi lần & \NRepeats{} lần lặp, bốn biến thể xen kẽ mỗi lần \\
Ghi môi trường & \texttt{bench/capture\_env.sh}, lưu cùng mỗi kết quả & \texttt{bench/capture\_env.sh}, lưu cùng mỗi kết quả \\
Bộ đếm phần cứng & Không có trong khách & Có: khoảng trống dụng cụ đo rõ nhất \\
\hline
Đủ cho & Tương đương, tiêm lỗi, cấu trúc FFI, điểm gối bão hoà và so sánh quá tải, và so sánh tài nguyên & Thêm đuôi xa dưới bão hoà, và cơ chế đằng sau kết quả tài nguyên và thao tác \\
\hline
\end{tabular}
\end{table}

\begin{table}[H]
\centering
\small
\caption[Cấu hình dịch vụ, khối lượng công việc và chiến dịch]{Cấu hình dịch vụ được kiểm thử, khối lượng công việc điều khiển nó, và chiến dịch
đo. Tham số dịch vụ cố định bởi đặc tả đóng băng, một ngoại lệ: số shard đặt khớp lõi máy cung cấp
(Mục~\ref{sec:testbed}). Tất cả, kể cả ngoại lệ đó, giống nhau ở hai ngôn ngữ. Số chiến dịch sinh từ
nhật ký đo cùng các
bảng báo cáo chúng. Tải đại diện không chọn bằng mắt mà suy ra, là tải quét bận nhất mà cả bốn biến
thể build vẫn giữ p99 dưới $5$~ms.\label{tab:params}}
\footnotesize
\renewcommand{\arraystretch}{1.05}
\begin{tabular}{@{}R{0.28\linewidth} R{0.28\linewidth} R{0.36\linewidth}@{}}
\hline\hline
Tham số & Giá trị & Ghi chú \\
\hline
\multicolumn{3}{@{}l}{\emph{Dịch vụ được kiểm thử, giống nhau ở hai ngôn ngữ}} \\
Shard và luồng worker & $7$ & một luồng sở hữu một shard, mỗi loại một trên lõi được ghim \\
Giới hạn hàng đợi mỗi shard & $1\,024$ yêu cầu & có giới hạn và chặn, nên quá tải quan sát được \\
Dung lượng bản đồ ban đầu & $4\,096$ mục & dự trữ để trạng thái ổn định không rehash \\
Bearer mỗi ngữ cảnh & tối đa $8$ & \\
Payload báo hiệu & $64$ byte & \\
Trần ngữ cảnh mỗi shard & $200\,000$ & \\
Chọn shard & SplitMix64, rồi modulo & ghim: khuếch tán định danh tuần tự \\
Hash trong shard & FNV-1a, 64-bit & ghim hai bên và kiểm bit-bit \\
\hline
\multicolumn{3}{@{}l}{\emph{Khối lượng công việc}} \\
Khối lượng & signal, mixed & mixed đi qua toàn vòng đời RAN \\
UE hoạt động & $5\,000$ & \\
Mức tải đề nghị & \NLoadLevels{} & tới $\MaxOfferedLoad$~kops/s \\
Tải đại diện & $\ReprLoad$~kops/s & suy ra, xem chú thích \\
Bão hoà đo được & $\SatKnee$~kops/s & theo tải đề nghị tới $\TrackingLoad$~kops/s \\
Mô hình tải & vòng hở, lịch dự định cố định & độ trễ đo từ thời điểm gửi dự định \\
\hline
\multicolumn{3}{@{}l}{\emph{Chiến dịch}} \\
Biến thể build & $\NVariants$ & hai ngôn ngữ $\times$ hai bộ cấp phát \\
Lần lặp mỗi điểm quét & \NRepeats{} & biến thể xen kẽ trong mỗi lần lặp \\
Lần chạy đo & $\NRuns$ & \\
Ghi độ trễ & HDR histogram, $[1\,\text{ns}, 3.6\,\text{s}]$, 3 chữ số có nghĩa & cùng thuật toán và cấu hình hai bên \\
Thống kê tóm tắt & trung vị với CI bootstrap $95\%$ & khoảng báo đầy đủ, không phải nửa độ rộng \\
So sánh ngôn ngữ & kiểm định hoán vị ghép cặp chính xác & \NTests{} kiểm định, hiệu chỉnh Holm \\
Sàn $p$ có thể đạt & $\MinAttainableP$ / $\CorrectedFloor$ & chưa hiệu chỉnh / sau hiệu chỉnh \\
\hline
\end{tabular}
\end{table}

\appendix{Ví dụ tiêm lỗi\label{appendix:code}}
Là minh hoạ cho nghiên cứu tiêm lỗi (Chương~\ref{chapter:eval}, RQ2), lỗi use-after-free được
trình bày ở cả hai ngôn ngữ. Trong C++ bí danh vào bản đồ ngữ cảnh
sống sót sau \texttt{erase} giải phóng nó, nên chương trình sau biên dịch sạch và chỉ được
AddressSanitizer đánh dấu lúc chạy:

\begin{lstlisting}[language=C++]
std::unordered_map<unsigned long, UeContext> ctxs;
ctxs[1001] = UeContext{1001, 0};
UeContext* p = &ctxs.at(1001);  // alias into the container
ctxs.erase(1001);               // frees the node that p refers to
p->messages_handled += 1;       // use-after-free
\end{lstlisting}

Chương trình Rust an toàn tương ứng không biên dịch: bộ kiểm tra mượn cấm sửa bản đồ
(\texttt{remove}) trong khi tham chiếu \texttt{p} vào nó vẫn sống, nên lỗi không thể viết ngay từ
đầu (lỗi \texttt{E0499}):

\begin{lstlisting}[language=Rust]
let mut ctxs: HashMap<u64, UeContext> = HashMap::new();
ctxs.insert(1001, UeContext { messages_handled: 0 });
let p = ctxs.get_mut(&1001).unwrap(); // mutable borrow starts
ctxs.remove(&1001);                  // borrow while p live -> rejected
p.messages_handled += 1;             // p used after the remove
\end{lstlisting}
